\section{Metric spaces}
\pagenumbering{arabic}
\subsection{Definitions and examples}

\begin{de}
A metric space is a pair $(X,d)$, where $X$ is a set and
$d:X\times X\to[0,\infty)$ is a function, called a metric, satisfying the
following properties for all $x,y,z\in X$:

(a) $d(x,y)=d(y,x)$;

(b) $d(x,y)=0$ if and only if $x=y$;

(c) $d(x,y)\leq d(x,z)+d(z,y)$ (the triangle inequality).

For $x\in X$ and $r>0$, write
\[
B(x;r)=\{y\in X:d(x,y)<r\},\qquad
\overline{B}(x;r)=\{y\in X:d(x,y)\leq r\}.
\]
The set $B(x;r)$ is the open ball centered at $x$ with radius $r$, and
$\overline{B}(x;r)$ is the closed ball centered at $x$ with radius $r$.
\end{de}

\begin{de}
Let $(X,d)$ be a metric space. A subset $G\subseteq X$ is open if, for every
$x\in G$, there is an $r>0$ such that $B(x;r)\subseteq G$. A subset
$F\subseteq X$ is closed if its complement $X\setminus F$ is open.
\end{de}

\begin{de}
If $(X,d)$ is a metric space and $Y\subseteq X$, then the restriction of $d$
to $Y\times Y$ is a metric on $Y$. A subset $A\subseteq Y$ is relatively open
(respectively, relatively closed) in $Y$ if it is open (respectively, closed)
in the metric space $Y$.
\end{de}

\begin{thm}
Let $(X,d)$ be a metric space and let $Y\subseteq X$.

(a) A subset $G\subseteq Y$ is relatively open in $Y$ if and only if there is
an open subset $U\subseteq X$ such that $G=U\cap Y$.

(b) A subset $F\subseteq Y$ is relatively closed in $Y$ if and only if there
is a closed subset $D\subseteq X$ such that $F=D\cap Y$.
\end{thm}

\begin{proof}
(a) Suppose first that $G$ is relatively open in $Y$. For each $x\in G$, choose
$r_x>0$ such that $B(x;r_x)\cap Y\subseteq G$, and set
\[
U=\bigcup_{x\in G}B(x;r_x).
\]
Then $U$ is open and $U\cap Y=G$. Conversely, if $G=U\cap Y$ with $U$ open in
$X$, then for each $x\in G$ there is an $r>0$ such that $B(x;r)\subseteq U$.
Hence $B(x;r)\cap Y\subseteq G$, so $G$ is relatively open in $Y$.

(b) The set $F$ is relatively closed in $Y$ if and only if $Y\setminus F$ is
relatively open in $Y$. By part (a), this is equivalent to
$Y\setminus F=U\cap Y$ for some open set $U\subseteq X$. In that case,
with $D=X\setminus U$, we have
\[
D\cap Y=(X\setminus U)\cap Y=Y\setminus(U\cap Y)=F.
\]
The converse follows by reversing this argument.
\end{proof}

\begin{thm}
Let $(X,d)$ be a metric space.

(a) For all $x,y,z\in X$,
\[
\bigl|d(x,y)-d(y,z)\bigr|\leq d(x,z).
\]

(b) A finite intersection of open sets is open.

(c) An arbitrary union of open sets is open.

(d) A finite union of closed sets is closed.

(e) An arbitrary intersection of closed sets is closed.
\end{thm}

\begin{proof}
For (a), the triangle inequality gives
$d(x,y)\leq d(x,z)+d(z,y)$ and
$d(y,z)\leq d(y,x)+d(x,z)$. Rearranging both inequalities gives the result.
Parts (b)--(e) follow directly from the definitions and De Morgan's laws.
\end{proof}

\begin{de}
Let $A\subseteq X$. The interior of $A$ is
\[
\operatorname{int}A=\bigcup\{G:G\text{ is open and }G\subseteq A\}.
\]
The closure of $A$ is
\[
\operatorname{cl}A=\bigcap\{F:F\text{ is closed and }A\subseteq F\}.
\]
The boundary of $A$ is
\[
\partial A=\operatorname{cl}A\cap\operatorname{cl}(X\setminus A).
\]
\end{de}

\begin{thm}
Let $A\subseteq X$.

(a) A point $x$ belongs to $\operatorname{int}A$ if and only if there is an
$r>0$ such that $B(x;r)\subseteq A$.

(b) A point $x$ belongs to $\operatorname{cl}A$ if and only if
$B(x;r)\cap A\neq\varnothing$ for every $r>0$.
\end{thm}

\begin{proof}
Part (a) follows immediately from the definition of an open set and the fact
that $\operatorname{int}A$ is the union of all open subsets of $A$.

For (b), if $B(x;r)\cap A=\varnothing$ for some $r>0$, then
$A\subseteq X\setminus B(x;r)$, where $X\setminus B(x;r)$ is closed and does
not contain $x$. Thus $x\notin\operatorname{cl}A$. Conversely, if
$x\notin\operatorname{cl}A$, then the open set
$X\setminus\operatorname{cl}A$ contains $x$, so some ball centered at $x$ is
disjoint from $A$.
\end{proof}

\begin{thm}
Let $A\subseteq X$.

(a) $A$ is closed if and only if $A=\operatorname{cl}A$.

(b) $A$ is open if and only if $A=\operatorname{int}A$.

(c)
\[
\operatorname{cl}A=X\setminus\operatorname{int}(X\setminus A),\qquad
\operatorname{int}A=X\setminus\operatorname{cl}(X\setminus A),
\]
and
\[
\partial A=\operatorname{cl}A\setminus\operatorname{int}A.
\]

(d) If $A_1,\ldots,A_n\subseteq X$, then
\[
\operatorname{cl}\left(\bigcup_{k=1}^n A_k\right)
=\bigcup_{k=1}^n\operatorname{cl}A_k.
\]

(e) If $A_1,\ldots,A_n\subseteq X$, then
\[
\operatorname{int}\left(\bigcap_{k=1}^n A_k\right)
=\bigcap_{k=1}^n\operatorname{int}A_k.
\]
\end{thm}

\begin{proof}
Parts (a) and (b) follow from the fact that $\operatorname{cl}A$ is the
smallest closed set containing $A$, whereas $\operatorname{int}A$ is the
largest open set contained in $A$.

For the first identity in (c), a point $x$ is outside $\operatorname{cl}A$ if
and only if some open neighborhood of $x$ is contained in $X\setminus A$;
this is equivalent to $x\in\operatorname{int}(X\setminus A)$. The second
identity follows by replacing $A$ by $X\setminus A$. Finally,
\[
\partial A
=\operatorname{cl}A\cap\operatorname{cl}(X\setminus A)
=\operatorname{cl}A\cap\bigl(X\setminus\operatorname{int}A\bigr)
=\operatorname{cl}A\setminus\operatorname{int}A.
\]

For (e), the inclusion from left to right is immediate. Conversely, if
$x\in\bigcap_{k=1}^n\operatorname{int}A_k$, choose $r_k>0$ such that
$B(x;r_k)\subseteq A_k$ and set $r=\min\{r_1,\ldots,r_n\}$. Then
$B(x;r)\subseteq\bigcap_{k=1}^n A_k$. Part (d) follows from (c), (e), and De
Morgan's laws.
\end{proof}

\begin{de}
If $(X_1,d_1)$ and $(X_2,d_2)$ are metric spaces, define a metric on
$X_1\times X_2$ by
\[
d\bigl((x_1,x_2),(y_1,y_2)\bigr)
=d_1(x_1,y_1)+d_2(x_2,y_2).
\]
\end{de}

\begin{thm}
Adopt the preceding notation.

(a) If $G_1\subseteq X_1$ and $G_2\subseteq X_2$ are open, then
$G_1\times G_2$ is open in $X_1\times X_2$.

(b) If $F_1\subseteq X_1$ and $F_2\subseteq X_2$ are closed, then
$F_1\times F_2$ is closed in $X_1\times X_2$.

(c) If $G$ is open in $X_1\times X_2$ and $(x_1,x_2)\in G$, then there is an
$r>0$ such that
\[
B_{d_1}(x_1;r)\times B_{d_2}(x_2;r)\subseteq G.
\]
\end{thm}

\begin{proof}
(a) If $(x_1,x_2)\in G_1\times G_2$, choose $\varepsilon_i>0$ such that
$B_{d_i}(x_i;\varepsilon_i)\subseteq G_i$ for $i=1,2$. Then the ball in the
product metric of radius $\min\{\varepsilon_1,\varepsilon_2\}$ is contained
in $G_1\times G_2$.

(b) We have
\[
(X_1\times X_2)\setminus(F_1\times F_2)
=((X_1\setminus F_1)\times X_2)\cup(X_1\times(X_2\setminus F_2)),
\]
which is open by part (a).

(c) Choose $\varepsilon>0$ such that
$B_d((x_1,x_2);\varepsilon)\subseteq G$, and put $r=\varepsilon/2$. If
$d_1(x_1,y_1)<r$ and $d_2(x_2,y_2)<r$, then
$d((x_1,x_2),(y_1,y_2))<\varepsilon$, proving the claim.
\end{proof}

\begin{de}
A subset $E$ of a metric space $(X,d)$ is dense if
$\operatorname{cl}E=X$. The metric space is separable if it has a countable
dense subset.
\end{de}

\begin{thm}
A set $E$ is dense in $(X,d)$ if and only if
$B(x;r)\cap E\neq\varnothing$ for every $x\in X$ and every $r>0$.
\end{thm}

\begin{proof}
This is the characterization of closure by open balls applied to the identity
$\operatorname{cl}E=X$.
\end{proof}

\begin{ex}
Let $\ell^\infty$ denote the set of all bounded sequences of real numbers;
that is,
\[
\ell^\infty=\left\{(x_n):\sup_n|x_n|<\infty\right\}.
\]
If $x=(x_n)$ and $y=(y_n)$ belong to $\ell^\infty$, define
\[
d(x,y)=\sup_n|x_n-y_n|.
\]

(a) Show that $d$ is a metric on $\ell^\infty$.

(b) If $e_n$ is the sequence with a $1$ in the $n$th coordinate and zeros
elsewhere, show that
$B(e_n;1/2)\cap B(e_m;1/2)=\varnothing$ whenever $n\neq m$.

(c) Is the set $\{e_n:n\geq1\}$ closed?
\end{ex}

\begin{proof}
(a) Symmetry and positive definiteness are immediate. For every coordinate,
\[
|x_n-y_n|\leq |x_n-z_n|+|z_n-y_n|.
\]
Taking suprema gives $d(x,y)\leq d(x,z)+d(z,y)$.

(b) If $x$ belonged to both balls and $n\neq m$, then
$|x_n-1|<1/2$ because $d(x,e_n)<1/2$, while $|x_n|<1/2$ because the $n$th
coordinate of $e_m$ is zero. This is impossible.

(c) Yes. If a sequence of points from $A=\{e_n:n\geq1\}$ converges, then it is
Cauchy. Since $d(e_n,e_m)=1$ whenever $n\neq m$, that sequence must be
eventually constant. Its limit therefore belongs to $A$, so $A$ is closed.
\end{proof}

\subsection{Sequences and completeness}

\begin{de}
A sequence $(x_n)$ in $X$ converges to $x$ if, for every $\varepsilon>0$,
there is an integer $N$ such that $d(x_n,x)<\varepsilon$ whenever $n\geq N$.
We write $x_n\to x$.

If $A\subseteq X$, a point $x\in X$ is a limit point of $A$ if every ball
centered at $x$ contains a point of $A$ different from $x$.

For a set $A\subseteq X$, the distance from $x$ to $A$ is
\[
\operatorname{dist}(x,A)=\inf\{d(x,a):a\in A\},
\]
where we use the convention $\inf\varnothing=+\infty$. When $A$ is
nonempty, it is possible that $\operatorname{dist}(x,A)=0$ even though
$x\notin A$.

A sequence $(x_n)$ is Cauchy if, for every $\varepsilon>0$, there is an
integer $N$ such that $d(x_n,x_m)<\varepsilon$ whenever $m,n\geq N$. A metric
space is complete if every Cauchy sequence converges.
\end{de}

\begin{thm}
Let $(X,d)$ be a metric space.

(1) A subset $F\subseteq X$ is closed if and only if every convergent sequence
of points of $F$ has its limit in $F$.

(2) Every subsequence of a sequence converging to $x$ also converges to $x$.

(3) Let $A\subseteq X$.

(a) A point $x$ is a limit point of $A$ if and only if there is a sequence of
distinct points of $A$ converging to $x$.

(b) The set $A$ is closed if and only if it contains all its limit points.

(c)
\[
\operatorname{cl}A=A\cup\{x:x\text{ is a limit point of }A\}.
\]

(4)
\[
\operatorname{cl}A=\{x\in X:\operatorname{dist}(x,A)=0\}.
\]

(5) If $(x_n)$ is Cauchy and a subsequence converges to $x$, then $x_n\to x$.

(6) The diameter of a nonempty set $E$ is
\[
\operatorname{diam}E=\sup\{d(x,y):x,y\in E\},
\]
where the value may be $+\infty$. Moreover,
$\operatorname{diam}E=\operatorname{diam}(\operatorname{cl}E)$.
\end{thm}

The proofs are the same as the standard proofs for subsets and sequences in
$\mathbb{R}$, with absolute value replaced by the metric.

\begin{thm}[Cantor's Intersection Theorem]
A metric space $(X,d)$ is complete if and only if every sequence $(F_n)$ of
nonempty subsets satisfying

(i) each $F_n$ is closed;

(ii) $F_1\supseteq F_2\supseteq\cdots$;

(iii) $\operatorname{diam}F_n\to0$

has an intersection $\bigcap_{n=1}^\infty F_n$ consisting of exactly one
point.
\end{thm}

\begin{proof}
Suppose first that $X$ is complete. Choose $x_n\in F_n$. If $m,n\geq N$, then
$x_m,x_n\in F_N$, so
$d(x_m,x_n)\leq\operatorname{diam}F_N$. Thus $(x_n)$ is Cauchy and converges
to some $x\in X$. For each fixed $N$, all $x_n$ with $n\geq N$ lie in the
closed set $F_N$, so $x\in F_N$. Hence $x\in\bigcap_nF_n$. If $x$ and $y$
both belong to the intersection, then
$d(x,y)\leq\operatorname{diam}F_n$ for every $n$, so $x=y$.

Conversely, let $(x_n)$ be Cauchy and set
\[
F_n=\operatorname{cl}\{x_m:m\geq n\}.
\]
The sets $F_n$ are nonempty, closed, and decreasing. The Cauchy property and
$\operatorname{diam}(\operatorname{cl}E)=\operatorname{diam}E$ imply
$\operatorname{diam}F_n\to0$. Let $\{x\}=\bigcap_nF_n$. Since
$x_n,x\in F_n$, we have
$d(x_n,x)\leq\operatorname{diam}F_n\to0$. Thus $x_n\to x$, and $X$ is
complete.
\end{proof}

\begin{thm}
Let $(X,d)$ be complete and let $Y\subseteq X$. Then $Y$, with the restricted
metric, is complete if and only if $Y$ is closed in $X$.
\end{thm}

\begin{proof}
If $Y$ is closed, every Cauchy sequence in $Y$ converges in $X$, and its limit
belongs to $Y$. Thus $Y$ is complete.

Conversely, suppose $Y$ is complete and $(y_n)$ is a sequence in $Y$ that
converges in $X$ to $y$. The sequence is Cauchy in $Y$, so it converges in
$Y$ to some $z\in Y$. Uniqueness of limits in $X$ gives $y=z\in Y$.
Therefore $Y$ is closed.
\end{proof}

\begin{de}
A subset $A$ of a metric space is bounded if $A=\varnothing$ or there are
$x_0\in X$ and $R>0$ such that $A\subseteq B(x_0;R)$.
\end{de}

\begin{thm}
(a) A subset $A$ of $(X,d)$ is bounded if and only if, for every $x\in X$,
there is an $r>0$ such that $A\subseteq B(x;r)$.

(b) A finite union of bounded sets is bounded.

(c) The range of every Cauchy sequence is bounded.
\end{thm}

\begin{proof}
These statements follow directly from the triangle inequality.
\end{proof}

\begin{ex}
If $A\subseteq X$, show that
\[
\operatorname{int}A
=\{x\in X:\operatorname{dist}(x,X\setminus A)>0\}.
\]
Give an analogous characterization of $\partial A$.
\end{ex}

\begin{proof}
A point $x$ is in $\operatorname{int}A$ if and only if some ball centered at
$x$ is contained in $A$, which is equivalent to
$\operatorname{dist}(x,X\setminus A)>0$. Also,
\[
\partial A
=\{x\in X:\operatorname{dist}(x,A)=0
\text{ and }\operatorname{dist}(x,X\setminus A)=0\}.
\]
\end{proof}

\begin{ex}
Let $(X,d)$ be the Cartesian product of $(X_1,d_1)$ and $(X_2,d_2)$ with the
sum metric.

(a) Show that $((x_n^1,x_n^2))$ is Cauchy in $X$ if and only if $(x_n^1)$ is
Cauchy in $X_1$ and $(x_n^2)$ is Cauchy in $X_2$.

(b) Show that $X$ is complete if and only if both $X_1$ and $X_2$ are
complete.
\end{ex}

\begin{proof}
Part (a) follows from
\[
d_i(x_n^i,x_m^i)\leq d((x_n^1,x_n^2),(x_m^1,x_m^2))
\]
and from the definition of the sum metric. Part (b) follows from part (a) and
the corresponding coordinatewise criterion for convergence.
\end{proof}

\subsection{Continuity}

\begin{de}
Let $(X,d)$ and $(Z,\rho)$ be metric spaces. A function $f:X\to Z$ is
continuous at $a\in X$ if, for every $\varepsilon>0$, there is a $\delta>0$
such that
\[
d(a,x)<\delta\quad\Longrightarrow\quad
\rho(f(a),f(x))<\varepsilon.
\]
The function is continuous on $X$ if it is continuous at every point of $X$.
\end{de}

\begin{thm}[Sequential criterion]
Let $(X,d)$ and $(Z,\rho)$ be metric spaces and let $f:X\to Z$. Then $f$ is
continuous at $a$ if and only if $x_n\to a$ implies
$f(x_n)\to f(a)$ for every sequence $(x_n)$ in $X$.
\end{thm}

\begin{thm}
Let $f:X\to Z$ be a function between metric spaces. The following are
equivalent:

(a) $f$ is continuous;

(b) $f^{-1}(U)$ is open in $X$ for every open set $U\subseteq Z$;

(c) $f^{-1}(D)$ is closed in $X$ for every closed set $D\subseteq Z$.
\end{thm}

\begin{proof}
This is the usual inverse-image characterization of continuity; (b) and (c)
are equivalent because inverse images commute with complements.
\end{proof}

\begin{thm}
If $A$ is a nonempty subset of the metric space $(X,d)$, then
\[
\bigl|\operatorname{dist}(x,A)-\operatorname{dist}(y,A)\bigr|
\leq d(x,y)
\]
for all $x,y\in X$.
\end{thm}

\begin{proof}
For every $a\in A$,
\[
d(x,a)\leq d(x,y)+d(y,a).
\]
Taking the infimum over $a\in A$ gives
\[
\operatorname{dist}(x,A)
\leq d(x,y)+\operatorname{dist}(y,A).
\]
Interchanging $x$ and $y$ yields the reverse inequality and hence the stated
estimate.
\end{proof}

\begin{cor}
For every nonempty $A\subseteq X$, the function
$f:X\to\mathbb{R}$ defined by $f(x)=\operatorname{dist}(x,A)$ is continuous.
In fact, it is $1$-Lipschitz.
\end{cor}

\begin{thm}
If $f,g:X\to\mathbb{R}$ are continuous, then $f+g$ and $fg$ are continuous.
If $f(x)\neq0$ for every $x\in X$, then $1/f$ is continuous.
\end{thm}

\begin{thm}
The composition of continuous functions is continuous.
\end{thm}

\begin{thm}[Urysohn's Lemma for metric spaces]
If $A$ and $B$ are disjoint closed subsets of a metric space $X$, then there
is a continuous function $f:X\to[0,1]$ such that $f=0$ on $A$ and $f=1$ on
$B$.
\end{thm}

\begin{proof}
If $A=\varnothing$, take $f\equiv1$; if $B=\varnothing$, take $f\equiv0$.
Otherwise, define
\[
f(x)=\frac{\operatorname{dist}(x,A)}
{\operatorname{dist}(x,A)+\operatorname{dist}(x,B)}.
\]
The denominator is positive at every point: if both distances were zero, the
point would belong to both closed sets. The preceding distance theorem and
the algebra of continuous functions show that $f$ is continuous. Its values
on $A$ and $B$ are as required.
\end{proof}

\begin{cor}
If $F$ is closed and $G$ is open with $F\subseteq G$, then there is a
continuous function $f:X\to[0,1]$ such that $f=1$ on $F$ and $f=0$ on
$X\setminus G$.
\end{cor}

\begin{proof}
Apply the preceding theorem to the disjoint closed sets $X\setminus G$ and
$F$.
\end{proof}

\begin{de}
A map $f:(X,d)\to(Z,\rho)$ is a homeomorphism if it is bijective and both
$f$ and $f^{-1}$ are continuous. Two metric spaces are homeomorphic if there
is a homeomorphism between them.

Two metrics $d$ and $\rho$ on the same set $X$ are equivalent if they induce
the same topology. For metric spaces, this is equivalent to saying that they
have the same convergent sequences, or that the identity map
$(X,d)\to(X,\rho)$ is a homeomorphism.
\end{de}

\begin{thm}
For every metric space $(X,d)$,
\[
\rho(x,y)=\frac{d(x,y)}{1+d(x,y)}
\]
defines an equivalent metric on $X$.
\end{thm}

\begin{proof}
Symmetry and positive definiteness are immediate. Let
$\phi(t)=t/(1+t)$. Since $\phi$ is increasing and
\[
\phi(a+b)\leq\phi(a)+\phi(b)\qquad(a,b\geq0),
\]
the triangle inequality for $d$ implies the triangle inequality for $\rho$.
Thus $\rho$ is a metric.

Moreover, $\rho(x,y)\leq d(x,y)$. Conversely, for every $\delta>0$,
\[
\rho(x,y)<\frac{\delta}{1+\delta}
\quad\Longrightarrow\quad d(x,y)<\delta.
\]
Hence the identity maps in both directions are continuous, so the metrics are
equivalent.
\end{proof}

\textbf{Remark.} The metric $\rho$ is bounded by $1$. Thus boundedness is not
a topological property.

\textbf{Remark.} Equivalent metrics preserve continuity but need not preserve
Cauchy sequences or completeness.

\begin{eg}
Let
\[
X=\{(x,y)\in\mathbb{R}^2:x^2+(y-1)^2=1\}
\]
with its Euclidean subspace metric. For $t\in\mathbb{R}$, take the line
through $(t,0)$ and $(0,2)$ and let $f(t)$ be its intersection with $X$ other
than $(0,2)$. Then $f$ is a homeomorphism from $\mathbb{R}$ onto
$Y=X\setminus\{(0,2)\}$. Define
\[
\rho(s,t)=d(f(s),f(t)).
\]
The metric $\rho$ is equivalent to the usual metric on $\mathbb{R}$, and
$f:(\mathbb{R},\rho)\to(Y,d)$ is an isometry. The sequence $(n)$ is Cauchy
with respect to $\rho$, because $f(n)\to(0,2)$ in $X$, but it has no limit in
$(\mathbb{R},\rho)$. Therefore $(\mathbb{R},\rho)$ is not complete.
\end{eg}

\begin{de}
A function $f:(X,d)\to(Z,\rho)$ is uniformly continuous if, for every
$\varepsilon>0$, there is a $\delta>0$ such that
\[
d(x,y)<\delta\quad\Longrightarrow\quad
\rho(f(x),f(y))<\varepsilon
\]
for all $x,y\in X$.
\end{de}

\begin{thm}
(a) A uniformly continuous function sends Cauchy sequences to Cauchy
sequences.

(b) If $A\subseteq X$, $Z$ is complete, and $f:A\to Z$ is uniformly
continuous, then $f$ has a unique uniformly continuous extension
$\overline f:\operatorname{cl}A\to Z$.
\end{thm}

\begin{proof}
Part (a) follows directly from the definitions. For (b), if
$x\in\operatorname{cl}A$, choose a sequence $(a_n)$ in $A$ with $a_n\to x$
and define
\[
\overline f(x)=\lim_{n\to\infty}f(a_n).
\]
Part (a) and completeness of $Z$ guarantee that the limit exists. Uniform
continuity shows that the value is independent of the chosen sequence and
that $\overline f$ is uniformly continuous. Density of $A$ gives uniqueness.
\end{proof}

\begin{ex}
(a) If $f:X\to Z$ is continuous, $A$ is dense in $X$, and $f(a)=z$ for every
$a\in A$, show that $f(x)=z$ for every $x\in X$.

(b) If $f:X\to Z$ is continuous and surjective and $A$ is dense in $X$, show
that $f(A)$ is dense in $Z$.
\end{ex}

\begin{proof}
(a) For any $x\in X$, choose a sequence $(a_n)$ in $A$ with $a_n\to x$.
Then $f(x)=\lim f(a_n)=z$.

(b) Let $z\in Z$ and let $V$ be an open neighborhood of $z$. Choose
$x\in X$ with $f(x)=z$. Then $f^{-1}(V)$ is an open neighborhood of $x$, so
it meets $A$. Hence $V$ meets $f(A)$.
\end{proof}

\subsection{Compactness}

\begin{de}
A collection $\mathcal{G}$ of subsets of $X$ is a cover of $E\subseteq X$ if
$E\subseteq\bigcup\mathcal{G}$. It is an open cover if every member of
$\mathcal{G}$ is open. A subcover is a subcollection that still covers $E$.

A subset $K$ of a metric space is compact if every open cover of $K$ has a
finite subcover.
\end{de}

\begin{thm}
Let $(X,d)$ be a metric space.

(a) Every compact subset of $X$ is closed and bounded.

(b) Every closed subset of a compact set is compact.

(c) The continuous image of a compact set is compact.
\end{thm}

\begin{cor}
If $X$ is nonempty and compact and $f:X\to\mathbb{R}$ is continuous,
then $f$ attains a minimum and a maximum on $X$.
\end{cor}

\begin{de}
A subset $K$ of a metric space is totally bounded if, for every $r>0$, it can
be covered by finitely many balls of radius $r$.

A collection $\mathcal{F}$ of sets has the finite intersection property (FIP)
if every finite subcollection has nonempty intersection.
\end{de}

\begin{thm}
For a subset $K$ of a metric space, the following are equivalent:

(a) $K$ is compact;

(b) every collection of closed subsets of $K$ having the FIP has nonempty
intersection;

(c) every sequence in $K$ has a convergent subsequence whose limit belongs to
$K$;

(d) every infinite subset of $K$ has a limit point in $K$;

(e) $K$, with its subspace metric, is complete and totally bounded.
\end{thm}

\begin{proof}
$(a)\Rightarrow(b)$: If a family $\mathcal{F}$ of closed subsets of $K$ has
empty intersection, then the relative-open sets $K\setminus F$, $F\in
\mathcal{F}$, cover $K$. A finite subcover would contradict the FIP.

$(a)\Rightarrow(d)$: If an infinite set $E\subseteq K$ had no limit point in
$K$, every point of $K$ would have a neighborhood meeting $E$ in at most one
point. A finite subcover of these neighborhoods would force $E$ to be finite.

$(d)\Rightarrow(c)$: If a sequence has finite range, it has a constant
subsequence. If its range is infinite, let $x$ be a limit point of the range.
Every neighborhood of $x$ contains infinitely many terms of the sequence, so
one can choose a subsequence converging to $x$.

$(c)\Rightarrow(e)$: A Cauchy sequence has a convergent subsequence; the whole
sequence then converges to the same limit, so $K$ is complete. If $K$ were not
totally bounded, there would be an $\varepsilon>0$ and a sequence in $K$ whose
distinct terms are at least $\varepsilon$ apart. Such a sequence has no
convergent subsequence, a contradiction.

$(e)\Rightarrow(c)$: Repeatedly cover $K$ by finitely many balls of radii
$1,1/2,1/3,\ldots$ and choose nested subsequences lying in one ball at each
stage. A diagonal subsequence is Cauchy and therefore converges in $K$.

Finally, $(c)\Rightarrow(a)$. First note the Lebesgue-number property: for
every open cover $\mathcal{G}$ of $K$, there is an $r>0$ such that every ball
$B(x;r)$ with $x\in K$ is contained in some member of $\mathcal{G}$. Otherwise,
choose $x_n\in K$ such that $B(x_n;1/n)$ is contained in no member of the
cover. A convergent subsequence tends to some $x\in K$, and an open member of
the cover containing $x$ eventually contains the corresponding balls, a
contradiction. Since $(c)$ implies total boundedness by the preceding
argument, finitely many balls of radius $r$ cover $K$, and the corresponding
members of $\mathcal{G}$ form a finite subcover.
\end{proof}

\textbf{Remark.} Compact sets are totally bounded, but total boundedness alone
does not imply compactness; completeness is also required.

\begin{thm}[Heine--Borel]
A subset of $\mathbb{R}^q$ is compact if and only if it is closed and bounded.
\end{thm}

\begin{thm}
Every continuous map from a compact metric space into a metric space is
uniformly continuous.
\end{thm}

\begin{thm}
Every compact metric space is separable.
\end{thm}

\begin{proof}
For each $n\geq1$, compactness gives a finite set $F_n\subseteq X$ such that
\[
X=\bigcup_{x\in F_n}B(x;1/n).
\]
The set $F=\bigcup_{n=1}^\infty F_n$ is countable. Given $x\in X$ and
$\varepsilon>0$, choose $n$ with $1/n<\varepsilon$ and then choose
$y\in F_n$ with $d(x,y)<1/n$. Thus $F$ is dense.
\end{proof}

\begin{ex}
Show that the closure of a totally bounded set is totally bounded.
\end{ex}

\begin{proof}
Let $E$ be totally bounded and let $r>0$. Choose
$x_1,\ldots,x_n\in E$ such that
\[
E\subseteq\bigcup_{i=1}^nB(x_i;r/2).
\]
Then
\[
\operatorname{cl}E
\subseteq\bigcup_{i=1}^n\overline B(x_i;r/2)
\subseteq\bigcup_{i=1}^nB(x_i;r),
\]
so $\operatorname{cl}E$ is totally bounded.
\end{proof}

\begin{ex}
If $(X,d)$ is complete and $E\subseteq X$, show that $E$ is totally bounded
if and only if $\operatorname{cl}E$ is compact.
\end{ex}

\begin{proof}
If $E$ is totally bounded, then $\operatorname{cl}E$ is totally bounded by the
preceding exercise. It is closed in the complete space $X$, so it is complete;
hence it is compact. Conversely, if $\operatorname{cl}E$ is compact, then it
is totally bounded, and every subset of a totally bounded set is totally
bounded.
\end{proof}

\subsection{Connectedness}

\begin{de}
A metric space $X$ is connected if its only subsets that are both open and
closed are $X$ and $\varnothing$. Equivalently, $X$ cannot be written as the
union of two disjoint nonempty open sets. A subset of a metric space is
connected when it is connected with its subspace metric.
\end{de}

\begin{thm}
A subset of $\mathbb{R}$ is connected if and only if it is an interval.
\end{thm}

\begin{thm}
The continuous image of a connected set is connected.
\end{thm}

\begin{proof}
If $f(X)=C\cup D$ were a separation of $f(X)$, then
$f^{-1}(C)$ and $f^{-1}(D)$ would form a separation of $X$.
\end{proof}

\begin{cor}[Intermediate Value Theorem]
If $X$ is connected, $f:X\to\mathbb{R}$ is continuous, and
$a,b\in f(X)$ with $a<b$, then every $c\in[a,b]$ belongs to $f(X)$.
\end{cor}

\begin{thm}
Let $X$ be a metric space.

(a) If $\{E_i:i\in I\}$ is a collection of connected subsets such that
$E_i\cap E_j\neq\varnothing$ for all $i,j\in I$, then
$\bigcup_{i\in I}E_i$ is connected.

(b) If $(E_n)$ is a sequence of connected subsets such that
$E_n\cap E_{n+1}\neq\varnothing$ for every $n$, then
$\bigcup_{n=1}^\infty E_n$ is connected.
\end{thm}

\begin{proof}
For (a), let $A$ be a nonempty relatively clopen subset of
$E=\bigcup_iE_i$. It meets some $E_{i_0}$, so connectedness gives
$E_{i_0}\subseteq A$. Every $E_j$ meets $E_{i_0}$, hence meets $A$, and
therefore $E_j\subseteq A$. Thus $A=E$.

For (b), the finite union $F_n=E_1\cup\cdots\cup E_n$ is connected by
induction. Since all the $F_n$ contain $E_1$, part (a) shows that
$\bigcup_nF_n=\bigcup_nE_n$ is connected.
\end{proof}

\begin{cor}
The union of two intersecting connected subsets is connected.
\end{cor}

\begin{de}
A component of a metric space is a maximal connected subset.
\end{de}

\textbf{Remark.} Maximal means maximal with respect to inclusion.

\begin{thm}
Every connected subset of a metric space is contained in a component;
distinct components are disjoint; and the union of all components is the
whole space.
\end{thm}

\begin{proof}
Let $D$ be connected. Take the union $C$ of all connected subsets that contain
$D$. Every two such sets meet in $D$, so $C$ is connected. It is maximal by
construction and hence is a component containing $D$. Applying this to the
singleton $\{x\}$ shows that every point lies in a component. If two
components meet, their union is connected, so maximality forces them to be
equal.
\end{proof}

\begin{thm}
If $C$ is connected and
$C\subseteq Y\subseteq\operatorname{cl}C$, then $Y$ is connected.
\end{thm}

\begin{proof}
If $Y=A\cup B$ were a separation, then connectedness of $C$ would imply that
$C$ is contained entirely in one of the two sets, say $A$. But $C$ is dense
in $Y$, so the nonempty relatively open set $B$ would have to meet $C$, a
contradiction.
\end{proof}

\begin{cor}
The closure of a connected set is connected, and every component is closed.
\end{cor}

\begin{de}
Let $E\subseteq X$, let $x,y\in E$, and let $\varepsilon>0$. An
$\varepsilon$-chain from $x$ to $y$ in $E$ is a finite sequence
$x=x_1,\ldots,x_n=y$ in $E$ such that

(i) $B(x_k;\varepsilon)\subseteq E$ for every $k$;

(ii) $d(x_{k-1},x_k)<\varepsilon$ for $2\leq k\leq n$.
\end{de}

\begin{thm}
Let $G$ be an open subset of $\mathbb{R}^q$.

(a) Every component of $G$ is open, and $G$ has at most countably many
components.

(b) The set $G$ is connected if and only if, for every $x,y\in G$, there is
an $\varepsilon>0$ and an $\varepsilon$-chain in $G$ from $x$ to $y$.
\end{thm}

\begin{proof}
(a) Let $H$ be a component of $G$ and $x\in H$. Since $G$ is open, some ball
$B(x;r)$ lies in $G$. The ball is connected and meets $H$, so
$B(x;r)\subseteq H$. Thus $H$ is open. Every nonempty open subset of
$\mathbb{R}^q$ contains a point of $\mathbb{Q}^q$, and distinct components
are disjoint; assigning a rational point to each component proves
countability.

(b) Suppose first that $G$ is connected. Fix $x\in G$ and let $D$ be the set
of points that can be joined to $x$ by a polygonal path in $G$. Small balls
show that both $D$ and $G\setminus D$ are open in $G$, so $D=G$. Let
$\gamma$ be a polygonal path from $x$ to $y$. Its compact image has positive
distance from $\mathbb{R}^q\setminus G$. Uniform continuity of $\gamma$
allows a sufficiently fine subdivision of the path whose vertices form an
$\varepsilon$-chain for some $\varepsilon>0$.

Conversely, the union of the balls appearing in an $\varepsilon$-chain is a
connected subset of $G$: consecutive balls intersect because the distance
between their centers is less than $\varepsilon$. Hence the endpoints of
every chain lie in the same component. If such a chain exists between every
pair of points, $G$ has only one component and is connected.
\end{proof}

\subsection{The Baire Category Theorem}

\begin{thm}[Baire Category Theorem]
If $(X,d)$ is complete and $(U_n)$ is a sequence of open dense subsets of
$X$, then $\bigcap_{n=1}^\infty U_n$ is dense in $X$.
\end{thm}

\begin{proof}
Let $G$ be a nonempty open subset of $X$. Since $U_1$ is dense,
$G\cap U_1$ is a nonempty open set. Choose $x_1\in X$ and $r_1<1$ such that
\[
\overline B(x_1;r_1)\subseteq G\cap U_1.
\]
Inductively, after choosing $x_{n-1}$ and $r_{n-1}$, the set
$B(x_{n-1};r_{n-1})\cap U_n$ is nonempty and open. Choose $x_n$ and
$0<r_n<1/n$ such that
\[
\overline B(x_n;r_n)
\subseteq B(x_{n-1};r_{n-1})\cap U_n.
\]
For $m>n$, the nesting gives $x_m\in B(x_n;r_n)$, so $(x_n)$ is Cauchy. Let
$x_n\to x$. For every fixed $n$, all sufficiently late terms lie in the
closed set $\overline B(x_n;r_n)$, hence
$x\in\overline B(x_n;r_n)$. Therefore
\[
x\in G\cap\bigcap_{n=1}^\infty U_n.
\]
Since every nonempty open $G$ meets the intersection, the intersection is
dense.
\end{proof}

\begin{cor}
If a nonempty complete metric space is a countable union
$X=\bigcup_{n=1}^\infty F_n$ of closed sets, then
$\operatorname{int}F_n\neq\varnothing$ for at least one $n$.
\end{cor}

\begin{proof}
If every $F_n$ had empty interior, then each open set
$U_n=X\setminus F_n$ would be dense. The Baire Category Theorem would imply
that $\bigcap_nU_n$ is dense, but this intersection is empty.
\end{proof}

\textbf{Remark.} The assumption that each $U_n$ is open is essential.

\begin{ex}
A set $E$ is nowhere dense if
$\operatorname{int}(\operatorname{cl}E)=\varnothing$. If
$A=\bigcup_{n=1}^\infty E_n$ and every $E_n$ is nowhere dense, show that
$X\setminus A$ is dense in every complete metric space $X$.
\end{ex}

\begin{proof}
For each $n$, the set
\[
U_n=X\setminus\operatorname{cl}E_n
\]
is open and dense. Hence $\bigcap_nU_n$ is dense by the Baire Category
Theorem. Since
\[
\bigcap_nU_n
=X\setminus\bigcup_n\operatorname{cl}E_n
\subseteq X\setminus\bigcup_nE_n=X\setminus A,
\]
the set $X\setminus A$ is dense as well.
\end{proof}
